\chapter*{Приложение A. Тексты промптов судьи}
\addcontentsline{toc}{chapter}{Приложение A}
\label{app:prompts}

Все три промпта реализуют одну задачу — попарную оценку ответов
системы $S$ и конкурента $C$ по четырём критериям чеклиста
(Relevance $\to$ Accuracy $\to$ Completeness $\to$ Readability).
Различие — в глубине инструкции и наличии примеров.

\section{Промпт c1\_compact (компактный, zero-shot)}
\label{app:compact}

Назначение: минимальная инструкция, судья выставляет
чекбоксы напрямую без явного рассуждения.

Стратегия генерации: zero-shot, temperature$=0{,}0$.

\begin{verbatim}
[SYSTEM]
Ты — эксперт по оценке качества ответов
диалоговых ассистентов.

[USER]
Запрос пользователя:
{query}

Ответ системы A:
{response_S}

Ответ системы B:
{response_C}

Оцени каждый ответ строго по порядку.
Для каждого критерия укажи true/false.
При первом false последующие критерии
автоматически считаются false.

Критерии (в порядке проверки):
1. Relevance: отвечает ли ответ на запрос в целом?
2. Accuracy: фактически верен ли ответ
   во всех модальностях?
3. Completeness: полностью ли покрыт запрос,
   с нужным форматом и деталями?
4. Readability: структурирован и читаем ли ответ?

Формат вывода (строго JSON):
{{
  "A": {{"relevance": bool, "accuracy": bool,
        "completeness": bool, "readability": bool}},
  "B": {{"relevance": bool, "accuracy": bool,
        "completeness": bool, "readability": bool}}
}}
\end{verbatim}

\section{Промпт c2\_stoprule (явное правило «стоп»)}
\label{app:stoprule}

Назначение: явная инструкция «стоп на первом false»
снижает неоднородность интерпретации чеклиста.

Стратегия генерации: zero-shot, temperature$=0{,}0$.

\begin{verbatim}
[SYSTEM]
Ты — эксперт по оценке качества ответов
диалоговых ассистентов.

[USER]
Запрос пользователя:
{query}

Ответ системы A:
{response_S}

Ответ системы B:
{response_C}

ВАЖНО: проверяй критерии строго по порядку.
Как только критерий получает false —
все последующие автоматически false.
Не оценивай длину ответа как показатель качества:
короткий точный ответ НЕ хуже длинного.

Критерии:
1. Relevance
2. Accuracy
3. Completeness
4. Readability

Формат вывода (строго JSON):
{{
  "A": {{"relevance": bool, "accuracy": bool,
        "completeness": bool, "readability": bool}},
  "B": {{"relevance": bool, "accuracy": bool,
        "completeness": bool, "readability": bool}}
}}
\end{verbatim}

\section{Промпт c3\_reason (цепочка рассуждений, few-shot)}
\label{app:reason}

Назначение: явное рассуждение по каждому критерию
перед вынесением вердикта (перенос идеи EvalPlanner).

Стратегия генерации: few-shot (1 пример), temperature$=0{,}0$.

\begin{verbatim}
[SYSTEM]
Ты — эксперт по оценке качества ответов
диалоговых ассистентов.

[USER — FEW-SHOT ПРИМЕР]
Запрос: «Покажи график цены акции за год»

Ответ A: «Вот график: [изображение]»
Ответ B: «Цена акции росла в течение года.»

Рассуждение:
- Relevance A: ответ содержит запрошенный график → true
- Relevance B: ответ описывает тренд словами,
  без графика → false; дальнейшие критерии false

Итог:
{{"A": {{"relevance": true, "accuracy": true,
       "completeness": true, "readability": true}},
 "B": {{"relevance": false, "accuracy": false,
       "completeness": false, "readability": false}}}}

[USER — РЕАЛЬНЫЙ ЗАПРОС]
Запрос пользователя:
{query}

Ответ системы A:
{response_S}

Ответ системы B:
{response_C}

Сначала напиши рассуждение по каждому критерию
для каждого ответа (1–2 предложения на критерий).
Затем выведи итоговый JSON в том же формате,
что в примере выше.
\end{verbatim}
